% Generated draft: requires mathematical and source review.
\documentclass[11pt]{article}
\usepackage[margin=25mm]{geometry}
\usepackage{fontspec}
\setmainfont{Latin Modern Roman}
\usepackage{amsmath,amssymb,mathtools,unicode-math}
\setmathfont{Latin Modern Math}
\usepackage{enumitem,needspace,xurl,hyperref}
\setlength{\parindent}{0pt}
\setlength{\parskip}{5pt}
\setlength{\emergencystretch}{3em}
\newcommand{\N}{\mathbb{N}}
\newcommand{\Z}{\mathbb{Z}}
\newcommand{\Q}{\mathbb{Q}}
\newcommand{\R}{\mathbb{R}}
\newcommand{\C}{\mathbb{C}}
\newcommand{\F}{\mathbb{F}}
\newcommand{\A}{\mathbb{A}}
\newcommand{\PP}{\mathbb{P}}
\newcommand{\E}{\mathbb{E}}
\newcommand{\eps}{\varepsilon}
\newcommand{\dd}{\,\mathrm{d}}
\begin{document}
\section*{UnsolvedMath ID: 700451}
\textbf{Source problem code:} catalogue entry
\subsection{Definitions and mathematical statement}
For a finitely generated group \(G\), define
\[
d(G)=\min\{|M|:M\subseteq G,\ \langle M\rangle=G\}.
\]
Following the source, a \emph{basis} of \(G\) is a generating set \(M\) with \(|M|=d(G)\); this terminology does not require \(G\) to be free. For a finitely presented group \(G\) and a basis \(M\), let \(F(M)\) be the free group on symbols indexed by \(M\), and let \(\pi_M:F(M)\to G\) send each symbol to its corresponding element. Define
\[
r_M(G)=\min\bigl\{|R|:R\subseteq\ker\pi_M,\ R\text{ finite},\ \langle\!\langle R\rangle\!\rangle_{F(M)}=\ker\pi_M\bigr\},
\qquad
r(G)=\min_{M\text{ a basis of }G}r_M(G),
\]
where \(\langle\!\langle R\rangle\!\rangle_{F(M)}\) denotes normal closure in \(F(M)\). Thus \(r_M(G)\) is the minimum number of defining relators when the generating elements are fixed to be \(M\).

\textbf{Part (a), the remaining question in the supplied record.} The source states the bound
\[
r_M(G)\leq d(G)+r(G)
\]
for every basis \(M\). Does there exist a finitely presented group \(G\) and a basis \(M\) of \(G\) such that
\[
r_M(G)=d(G)+r(G)?
\]
The source does not require \(G\) to be nontrivial or impose \(d(G)>0\). These omissions are retained; whether an additional restriction was intended requires review.

\textbf{Archived parts (b) and (c).} The source introduces nontrivial groups \(G_1,G_2\), and asks:
\begin{enumerate}
\item[(b)] For every such pair and every choice of bases \(M_1\) of \(G_1\) and \(M_2\) of \(G_2\), is
\[
r_{M_1\cup M_2}(G_1*G_2)=r_{M_1}(G_1)+r_{M_2}(G_2)?
\]
Here \(G_1*G_2\) is the free product and the union uses the images of the generating sets in the respective free factors.
\item[(c)] For every such pair, is
\[
r(G_1*G_2)=r(G_1)+r(G_2)?
\]
\end{enumerate}
The archived introduction defines these invariants for finitely presented groups, but its subsequent phrase ``any non-trivial groups'' does not explicitly impose finite presentability, or even finite generation, on the factors. Consequently, the displayed archived questions have an unresolved domain qualification. They can be read using the finite-valued definitions above when both factors are finitely presented; the source does not specify an extension to arbitrary groups.

The supplied archive reports negative answers to both (b) and (c). The supplied registry labels the overall problem partially solved and places (a) among the remaining parts. These are supplied source assessments, not independently verified mathematical or literature conclusions.
\subsection{Short English statement}
The remaining question asks whether a finitely presented group has a generating set of minimum size whose minimum number of defining relators attains the supplied upper bound \(d(G)+r(G)\). The record also preserves archived questions about additivity of these relator counts under free products, reporting negative answers to both. The intended exclusion of degenerate cases in (a), and the precise domain of the archived questions, require human review.
\subsection{Sources}
\begin{itemize}
\item Local UnsolvedMath registry, supplied complete record: permanent ID \(700451\), exact source problem code \texttt{catalogue entry}, title \emph{Kourovka Notebook Problem 7.57}; upstream revision \texttt{d371ee917ac9647fbb06cd6fc529deec6f3dda82}. The registry supplies the status \texttt{partially\_solved}.
\item V. A. Churkin, proposer, 1980. The supplied attribution identifies \emph{The Kourovka Notebook (Unsolved Problems in Group Theory)}, 7th issue (1980), Problem 7.57, updated in the 21st issue (2026), identified as \texttt{arXiv:1401.0300v48}. The supplied collection description names E. I. Khukhro and V. D. Mazurov as editors.
\item The supplied Archive of Solved Problems entry attributes negative answers to (b) and (c) to C. Hog, M. Lusztig, and W. Metzler, in \emph{Presentation classes, 3-manifolds and free products}, Lecture Notes in Mathematics, vol. 1167, Springer, Berlin, 1985, pp. 154--167. Author spelling and bibliographic scope are retained as supplied; no individual contribution title is provided.
\item The supplied status review is dated \texttt{2026-10-16}, records evidence as ``published'', and states that (b) and (c) are archived as solved while the remaining part is in the main section. This assessment is supplied dated curation, not independently verified. No external source was consulted for this draft.
\end{itemize}
\end{document}
